\begin{center}
	\large{\textbf{TINJAUAN PERFORMA MESIN PANAS KUANTUM RUANG FRAKSIONAL BERBASIS POTENSIAL TAK HINGGA 1 DIMENSI}}
\end{center}
%\begin{flushleft}
\begin{minipage}{5cm}
\begin{align*}
&\text{Nama} \qquad &&: \text{MUHAMMAD SULTHAN ZACKY RAIHAN PUTERA HASAN} \\
&\text{NRP} \qquad &&: \text{5001201046} \\
&\text{Departemen} \qquad &&: \text{Fisika} \\
&\text{Pembimbing} \qquad &&: \text{Dr. Heru Sukamto, M.Si} \\
\end{align*}
\end{minipage} \\
\textbf{Abstrak} 
\par Persamaan Schrödinger dapat didekati dengan beberapa cara, salah satunya menggunakan kalkulus fraksional yang memperkenalkan parameter baru yang disebut parameter fraksional \(\alpha\). Persamaan energi eigen yang diperoleh dengan menyelesaikan persamaan Schrödinger sekarang juga dipengaruhi oleh parameter baru ini. Persamaan energi eigen ini kemudian diterapkan pada salah satu cabang mekanika kuantum, mesin panas kuantum, untuk menyelidiki bagaimana parameter fraksional memengaruhi efisiensi dan daya keluaran mesin tersebut. Studi ini meninjau tiga siklus: Carnot, Otto, dan Brayton, menggunakan sebuah partikel yang terperangkap dalam sumur potensial tak terhingga satu dimensi sebagai zat kerja. Ditemukan bahwa efisiensi dapat direduksi hanya menjadi fungsi dari rasio kompresi \(r\) dan parameter fraksional \(\alpha\). Efisiensi dari semua siklus juga ditemukan memiliki bentuk yang sama, sementara daya keluarannya bervariasi di antara mereka. Tampaknya, parameter fraksional \(\alpha\) dapat secara signifikan memengaruhi baik efisiensi maupun daya keluaran dari mesin panas kuantum.\\
\begin{minipage}{1cm}
	\begin{align*}
	\text{\textbf{Kata kunci}}:\hspace{2mm}&\text{\textbf{\textit{Efisiensi, Mesin Panas Kuantum, Parameter Fraksional}}}
	\end{align*}
\end{minipage}
%\end{flushleft}

\newpage
\thispagestyle{empty}
\vspace*{\fill}
    \begin{center}
	Halaman ini sengaja dikosongkan
    \end{center}
\vspace*{\fill}
\pagebreak